\rhead{SE-EL1: DevOps \& Site Reliability Engineering}
\phantomsection 
\section*{DevOps \& Site Reliability Engineering (SE-EL1)}
\label{elec:SE_1}

\noindent \small{\hyperref[main:scheme]{$\leftarrow$ Back to Scheme}} \vspace{0.5cm}

\noindent \textbf{Credits:} 4 (3-0-2)

\subsection*{Theory}
\noindent\textbf{Unit 1} \\
\textit{DevOps Philosophy and CI/CD Pipelines:} DevOps cultural principles (CAMS - Culture, Automation, Measurement, Sharing), CALMS framework extension, Continuous Integration practices (trunk-based development, feature flags), Continuous Delivery/Deployment pipelines, GitOps declarative deployments, Pipeline as Code (Jenkinsfile, GitHub Actions, GitLab CI), Artifact management and immutable infrastructure.

\vspace{0.3cm}

\noindent\textbf{Unit 2} \\
\textit{Infrastructure as Code and Configuration Management:} Terraform workflow (init, plan, apply, destroy), HCL syntax and providers, State management (remote backends, locking), Ansible for configuration management (playbooks, roles, inventories), Puppet/Chef declarative models, Kubernetes as Infrastructure-as-Code (Helm charts, Kustomize), Desired State Configuration (DSC).

\vspace{0.3cm}

\noindent\textbf{Unit 3} \\
\textit{Site Reliability Engineering Principles:} SRE golden signals (latency, traffic, errors, saturation), Service Level Indicators/Objectives/Agreements (SLI/SLO/SLA), Error budgets and toil reduction, Capacity planning and load testing, Incident management (postmortem culture, blameless retrospectives), On-call engineering and incident response playbooks, Reliability engineering vs. traditional ops.

\vspace{0.3cm}

\noindent\textbf{Unit 4} \\
\textit{Observability and Monitoring Stack:} The three pillars (metrics, logs, traces), Prometheus time-series monitoring (federation, alerting), Grafana dashboards and anomaly detection, Loki for log aggregation, Jaeger/Zipkin distributed tracing, OpenTelemetry instrumentation, Synthetic monitoring and chaos engineering (Chaos Mesh, LitmusChaos), SLO-based alerting.

\vspace{0.3cm}

\noindent\textbf{Unit 5} \\
\textit{Production Excellence and Chaos Engineering:} Progressive delivery (feature flags, canary releases, blue-green deployments), A/B testing infrastructure, Multi-region failover and disaster recovery, GitOps with ArgoCD/Flux, Chaos engineering experiments (fault injection, resilience testing), Production readiness reviews (PRR), FinOps and cost optimization at scale.

\newpage

\subsection*{Practical}
\noindent\textbf{Lab 1: Complete CI/CD Pipeline Implementation} \\
Focuses on GitHub Actions/GitLab CI pipeline from code commit to production deployment.

\vspace{0.2cm}

\noindent\textbf{Lab 2: Terraform Multi-Cloud Infrastructure} \\
Focuses on declarative infrastructure provisioning across AWS/Azure/GCP with remote state.

\vspace{0.2cm}

\noindent\textbf{Lab 3: Ansible Configuration Management} \\
Focuses on infrastructure provisioning, application deployment, and compliance checking.

\vspace{0.2cm}

\noindent\textbf{Lab 4: Kubernetes GitOps with ArgoCD} \\
Focuses on declarative deployments, rollouts, and automated synchronization.

\vspace{0.2cm}

\noindent\textbf{Lab 5: Golden Signals Monitoring Stack} \\
Focuses on Prometheus/Grafana deployment with custom exporters and alerting rules.

\vspace{0.2cm}

\noindent\textbf{Lab 6: Distributed Tracing Implementation} \\
Focuses on OpenTelemetry instrumentation across microservices with Jaeger analysis.

\vspace{0.2cm}

\noindent\textbf{Lab 7: SLO/SLI Calculation and Error Budgets} \\
Focuses on service level definition, monitoring, and reliability budgeting.

\vspace{0.2cm}

\noindent\textbf{Lab 8: Chaos Engineering Experiments} \\
Focuses on fault injection (pod kills, network latency, resource exhaustion) and recovery analysis.

\vspace{0.2cm}

\noindent\textbf{Lab 9: Progressive Delivery Implementation} \\
Focuses on canary releases, feature flags, and A/B testing infrastructure.

\vspace{0.2cm}

\noindent\textbf{Lab 10: Production SRE Platform} \\
Focuses on complete SRE platform with monitoring, alerting, GitOps, and chaos engineering.

\vspace{0.2cm}

\newpage
\rhead{Curriculum Schema}
